% Part: applied-modal-logics
% Chapter: temporal-logic
% Section: language-epistemic-logic

\documentclass[../../../include/open-logic-section]{subfiles}

\begin{document}

\olfileid{aml}{tl}{sem}

\olsection{కాలిక తర్క అర్థవిజ్ఞానం}

\begin{defn}
కాలిక తర్కపు ప్రాథమిక భాషలో ఇవి ఉంటాయి:
\begin{enumerate}
  \tagitem{prvFalse}{\tetoken{అసత్యం}{!!{falsity}} కోసం ప్రతిజ్ఞావాక్య స్థిరాంకం~$\lfalse$.}{}
  \tagitem{prvTrue}{\tetoken{సత్యం}{!!{truth}} కోసం ప్రతిజ్ఞావాక్య స్థిరాంకం~$\ltrue$.}{}
  \item \tetoken{ప్రతిజ్ఞావాక్య చరాల}{!!{propositional variable}s}
    \tetoken{అనంతంగా లెక్కించదగిన}{!!{denumerable}s} సమితి: $\Obj
    p_0$, $\Obj p_1$, $\Obj p_2$, \dots
  \item ప్రతిజ్ఞావాక్య సంధాయకాలు: \startycommalist
  \iftag{prvNot}{\ycomma $\lnot$ (నిరాకరణ)}{}%
  \iftag{prvAnd}{\ycomma $\land$ (సంయోజనం)}{}%
  \iftag{prvOr}{\ycomma $\lor$ (విడియోజనం)}{}%
  \iftag{prvIf}{\ycomma $\lif$ (\tetoken{సోపాధికం}{!!{conditional}})}{}%
  \iftag{prvIff}{\ycomma $\liff$ (\tetoken{ద్విసోపాధికం}{!!{biconditional}})}{}.
  \item గతకాల కారకాలు $\Ptemp$, $\Htemp$.
  \item భవిష్యత్ కారకాలు $\Ftemp$, $\Gtemp$.
\end{enumerate}
\end{defn}

ఇతర రకాల మోడల్ కారకాలను చేర్చే అవకాశాన్ని తరువాత చర్చిస్తాం.

\begin{defn}
కాలిక భాష \emph{\tetoken{సూత్రాలను}{!!^{formula}s}} ఇలా
ఆగమన పద్ధతిలో నిర్వచిస్తాం:
\begin{enumerate}
\tagitem{prvFalse}{$\lfalse$ ఒక అణు \tetoken{సూత్రం}{!!{formula}}.}{}

\tagitem{prvTrue}{$\ltrue$ ఒక అణు \tetoken{సూత్రం}{!!{formula}}.}{}

\item ప్రతి ప్రతిజ్ఞావాక్య చరం $\Obj p_i$ ఒక (అణు)
  \tetoken{సూత్రం}{!!{formula}}.

\tagitem{prvNot}{$!A$ \tetoken{ఒక సూత్రం}{!!a{formula}} అయితే,
  $\lnot !A$ కూడా \tetoken{ఒక సూత్రం}{!!a{formula}}.}{}

\tagitem{prvAnd}{$!A$, $!B$ \tetoken{సూత్రాలు}{!!{formula}s}
  అయితే, $(!A \land !B)$ \tetoken{ఒక సూత్రం}{!!a{formula}}.}{}

\tagitem{prvOr}{$!A$, $!B$ \tetoken{సూత్రాలు}{!!{formula}s}
  అయితే, $(!A \lor !B)$ \tetoken{ఒక సూత్రం}{!!a{formula}}.}{}

\tagitem{prvIf}{$!A$, $!B$ \tetoken{సూత్రాలు}{!!{formula}s}
  అయితే, $(!A \lif !B)$ \tetoken{ఒక సూత్రం}{!!a{formula}}.}{}

\tagitem{prvIff}{$!A$, $!B$ \tetoken{సూత్రాలు}{!!{formula}s}
  అయితే, $(!A \liff !B)$ \tetoken{ఒక సూత్రం}{!!a{formula}}.}{}

\item $!A$ \tetoken{ఒక సూత్రం}{!!a{formula}} అయితే,
  $\Ptemp !A$, $\Htemp !A$, $\Ftemp !A$, $\Gtemp !A$
  అన్నీ \tetoken{సూత్రాలు}{!!{formula}s}.
  \sourcecorrection{OLTEAMLTLSEM-001}{మూలం భవిష్యత్
    సంభావ్యత కారకాన్ని ఇక్కడ $F$గా వ్రాసింది; పైన
    పరిచయమైన మాక్రో, దిగువ సత్య షరతు ప్రకారం
    $\Ftemp$గా సరిచేశాం.}

\tagitem{limitClause}{ఇవి తప్ప మరేదీ
  \tetoken{సూత్రం}{!!a{formula}} కాదు.}{}
\end{enumerate}
\end{defn}

ఇతర అంతర్గతార్థ తర్కాల మాదిరిగానే, కాలిక తర్కాల
అర్థవిజ్ఞానాన్ని సంబంధ నమూనాల ద్వారా ఇస్తాం.

\begin{defn}
  కాలిక భాషకు \emph{నమూనా} ఒక త్రయం
  $\mModel{M} = \tuple{T, \prec, V}$; ఇందులో
  \begin{enumerate}
  \item $T$ శూన్యం కాని సమితి; దీని సభ్యులను కాల బిందువులుగా
    అర్థం చేసుకుంటాం.
  \item $\prec$ అనేది $T$పై ద్విస్థాన సంబంధం.
  \item $V$ అనేది ప్రతి \tetoken{ప్రతిజ్ఞావాక్య
    చరానికి}{!!{propositional
    variable}}~$p$ కాల బిందువుల
    సమితి $V(p)$ను కేటాయించే ప్రమేయకం.
  \end{enumerate}
  $t \prec t'$ అయితే, $t$ అనేది~$t'$కు
  \emph{ముందు వస్తుంది} అంటాం.
  $t \in V(p)$ అయితే, $p$ అనేది~$t$ వద్ద
  \emph{సత్యం} అంటాం.
\end{defn}

ప్రస్తుతం పూర్వగామిత్వ సంబంధం~$\prec$పై ఎలాంటి
షరతులూ విధించలేదు. కాబట్టి మన కాలిక నమూనాల నిర్మాణంపై
ఇప్పటికి పరిమితి లేదు: కాలం రేఖీయంగా, శాఖలుగా,
చక్రీయంగా లేదా మరే నిర్మాణంతోనైనా ఉండవచ్చు.

సాధారణ మోడల్ తర్కంలాగే, ప్రతి కాలిక నమూనా ఏ
\tetoken{సూత్రాలు}{!!{formula}s} ఏ బిందువుల వద్ద
సత్యమో నిర్ణయిస్తుంది. సంబంధ అర్థవిజ్ఞానపు ప్రాథమిక
భావన కోసం ``నమూనా~$\mModel{M}$లో
\tetoken{సూత్రం}{!!{formula}}~$!A$ బిందువు~$t$ వద్ద
సత్యం'' అనే సంకేతాన్ని వాడతాం. ఆ సంబంధాన్ని ఆగమన
పద్ధతిలో నిర్వచిస్తాం; మోడల్ కాని కారకాలన్నింటికీ
అది సాధారణ మోడల్ తర్కంలోని నిర్వచనమే.

\begin{defn}\ollabel{defn:tmodels}
  నమూనా~$\mModel M$లో \emph{$t$ వద్ద
  \tetoken{ఒక సూత్రం}{!!a{formula}}~$!A$ సత్యం},
  సంకేతంగా $\mSat{M}{!A}[t]$, ఇలా ఆగమన పద్ధతిలో
  నిర్వచించబడుతుంది:
  \begin{enumerate}
  \tagitem{prvFalse}{\indcase{!A}{\lfalse}{$\mSat{M}{\lfalse}[t]$ ఎప్పుడూ కాదు}.}{}
  \tagitem{prvTrue}{\indcase{!A}{\ltrue}{$\mSat{M}{\ltrue}[t]$ ఎల్లప్పుడూ}.}{}
  \item $\mSat{M}{p}[t]$ అప్పుడే, అప్పుడు మాత్రమే $t \in V(p)$.
  \tagitem{prvNot}{\indcase{!A}{\lnot !B}{$\mSat{M}{\indfrm}[t]$
    అప్పుడే, అప్పుడు మాత్రమే $\mSat/{M}{!B}[t]$}.}{}
  \tagitem{prvAnd}{\indcase{!A}{(!B \land !C)}{$\mSat{M}{\indfrm}[t]$
    అప్పుడే, అప్పుడు మాత్రమే $\mSat{M}{!B}[t]$ మరియు
    $\mSat{M}{!C}[t]$}.}{}
  \tagitem{prvOr}{\indcase{!A}{(!B \lor !C)}{$\mSat{M}{\indfrm}[t]$
    అప్పుడే, అప్పుడు మాత్రమే $\mSat{M}{!B}[t]$ లేదా
    $\mSat{M}{!C}[t]$} (లేదా రెండూ).}{}
  \tagitem{prvIf}{\indcase{!A}{(!B \lif !C)}{$\mSat{M}{\indfrm}[t]$
    అప్పుడే, అప్పుడు మాత్రమే $\mSat/{M}{!B}[t]$ లేదా
    $\mSat{M}{!C}[t]$}.}{}
  \tagitem{prvIff}{\indcase{!A}{(!B \liff !C)}{$\mSat{M}{\indfrm}[t]$
    అప్పుడే, అప్పుడు మాత్రమే $\mSat{M}{!B}[t]$, $\mSat{M}{!C}[t]$
    రెండూ లేదా $\mSat{M}{!B}[t]$, $\mSat{M}{!C}[t]$
    రెండూ అసత్యమైనప్పుడు}.}{}
  \item\ollabel{defn:sub:mmodels-p}
    \indcase{!A}{\Ptemp  !B}{$\mSat{M}{\indfrm}[t]$
    అప్పుడే, అప్పుడు మాత్రమే $t' \prec t$ అయ్యే ఏదో
    $t' \in T$కు $\mSat{M}{!B}[t']$}
\item\ollabel{defn:sub:mmodels-h}
    \indcase{!A}{\Htemp  !B}{$\mSat{M}{\indfrm}[t]$
    అప్పుడే, అప్పుడు మాత్రమే $t' \prec t$ అయ్యే ప్రతి
    $t' \in T$కు $\mSat{M}{!B}[t']$}
   \item\ollabel{defn:sub:mmodels-f}
    \indcase{!A}{\Ftemp  !B}{$\mSat{M}{\indfrm}[t]$
    అప్పుడే, అప్పుడు మాత్రమే $t \prec t'$ అయ్యే ఏదో
    $t' \in T$కు $\mSat{M}{!B}[t']$}
\item\ollabel{defn:sub:mmodels-g}
    \indcase{!A}{\Gtemp  !B}{$\mSat{M}{\indfrm}[t]$
    అప్పుడే, అప్పుడు మాత్రమే $t \prec t'$ అయ్యే ప్రతి
    $t' \in T$కు $\mSat{M}{!B}[t']$}
  \end{enumerate}
\end{defn}

అర్థవిజ్ఞానం ప్రకారం $\Ptemp$,~$\Htemp$ పరస్పర
ద్వంద్వ కారకాలనీ, అలాగే $\Ftemp$, $\Gtemp$ కూడా
ద్వంద్వాలనీ గమనించవచ్చు. అందువల్ల $\Htemp !A$ను
$\lnot \Ptemp \lnot !A$గా నిర్వచించవచ్చు;
$\Gtemp$,~$\Ftemp$కూ ఇదే వర్తిస్తుంది.

\end{document}
